\begin{table}[htbp]
\centering
\caption{Replication Attempt: Washington's Border with Oregon}
\label{tab:wa_or}
\small
\begin{tabular}{lcccc}
\toprule
Specification & Coef. & SE (HC1) & SE (cl.) & $p$ (cl.) \\
\midrule
Transition year only (FY2016 vs.\ FY2015) & $-0.545$ & $0.329$ & $0.336$ & $0.105$ \\
DiD, post $=$ FY2016 onward & $+0.136$ & $0.381$ & $0.432$ & $0.753$ \\
DiD, FY2016 dropped, post $=$ FY2017 onward & $+0.161$ & $0.380$ & $0.443$ & $0.716$ \\
\bottomrule
\end{tabular}
\begin{minipage}{0.95\textwidth}
\vspace{0.3em}
\footnotesize
\textit{Notes:} Outcome is log(revenue per active dispensary). Washington county-year panel, fiscal years
2015--2025, $n=395$ (2015 is the only fully pre-legalization year). Treatment is the 13 Washington counties
within 70 miles of Oregon; controls include the 4 Idaho-border counties, the remaining interior counties,
log(population), and fiscal-year fixed effects. The transition-year estimate is directionally consistent with
\citet{HansenMillerWeber2020} but only marginally significant; the longer-run DiD estimates are positive and
insignificant, the opposite of sustained demand suppression. See Appendix~\ref{app:wa_or} for discussion.
\end{minipage}
\end{table}
